\chapter{Was sind Talente?}

Talente verleihen einem Helden besondere Fähigkeiten, die über seine Attribute und Fertigkeiten hinausgehen. Sie machen jeden Helden einzigartig und ermöglichen es ihm, sich im Laufe seiner Abenteuer weiterzuentwickeln.

Es gibt zwei Arten von Talenten:

\begin{itemize}
  \item Rollentalente
  \item allgemeine Talente
\end{itemize}

Während Rollentalente bestimmen, welche Aufgabe ein Held innerhalb seiner Gruppe übernimmt, verleihen allgemeine Talente ihm besondere Fähigkeiten, Eigenschaften oder Kenntnisse.

Die vollständige Übersicht aller Talente befindet sich in Anhang 3 -- Talente.

\section{Was ist ein Talent?}

Ein Talent ist eine besondere Fähigkeit eines Helden.

Talente werden durch Talentkarten dargestellt. Auf jeder Talentkarte steht, welche Wirkung das Talent besitzt und unter welchen Bedingungen es eingesetzt werden kann.

Talente können einem Helden beispielsweise erlauben,

\begin{itemize}
  \item besondere Angriffe auszuführen,
  \item Schaden zu verhindern,
  \item Verbündete zu unterstützen,
  \item Gegner zu behindern,
  \item Zauber zu wirken oder zu verbessern,
  \item außergewöhnliche Fähigkeiten einzusetzen oder
  \item außerhalb eines Kampfes neue Möglichkeiten zu erhalten.
\end{itemize}

Jedes Talent gehört entweder zu den Rollentalenten oder zu den allgemeinen Talenten.

\section{Welche Arten von Talenten gibt es?}

In Heldenpfade gibt es zwei unterschiedliche Talentarten.

Rollentalente bestimmen, welche Aufgabe ein Held innerhalb seiner Gruppe übernimmt. Sie sind vor allem auf den taktischen Teil des Spiels ausgerichtet.

Allgemeine Talente verleihen einem Helden besondere Fähigkeiten, Eigenschaften oder Kenntnisse. Sie beschreiben den Helden selbst und sind nicht an eine bestimmte Rolle gebunden.

Ein Held kann Talente beider Arten miteinander kombinieren.

\section{Was sind Rollentalente?}

Rollentalente bestimmen, auf welche Weise ein Held seiner Gruppe besonders hilft.

Sie sind in vier Rollen unterteilt:

\begin{itemize}
  \item \textbf{Draufgänger} gehen entschlossen in die Offensive. Ihre Talente helfen dabei, viel Schaden zu verursachen und Gegner schnell auszuschalten.
  \item \textbf{Verteidiger} schützen sich selbst und ihre Verbündeten. Sie verhindern Schaden, halten Gegner auf und schaffen Freiräume für die Gruppe.
  \item \textbf{Anführer} stärken ihre Verbündeten und unterstützen die gesamte Gruppe. Sie ermöglichen zusätzliche Handlungen, helfen bei der Heilung oder verbessern die Fähigkeiten anderer Helden.
  \item \textbf{Taktiker} beeinflussen das Kampfgeschehen. Sie behindern Gegner, verändern Positionen und schaffen günstige Situationen für ihre Verbündeten.
\end{itemize}

Eine Rolle ist kein fester Beruf und keine starre Klasse.

Ein Held muss sich nicht dauerhaft für eine einzige Rolle entscheiden. Er kann Rollentalente aus verschiedenen Rollen erlernen und miteinander kombinieren. Dadurch entstehen sehr unterschiedliche Spielstile.

\section{Was sind allgemeine Talente?}

Allgemeine Talente verleihen einem Helden besondere Fähigkeiten, Eigenschaften oder Kenntnisse.

Sie sind nicht an eine bestimmte Rolle gebunden und beschreiben, was einen Helden besonders macht.

Allgemeine Talente können sowohl außerhalb als auch innerhalb eines Kampfes Auswirkungen besitzen. Sie dienen jedoch nicht dazu, eine bestimmte Rolle innerhalb der Gruppe zu erfüllen.

Allgemeine Talente können beispielsweise:

\begin{itemize}
  \item den Umgang mit Magie ermöglichen oder verbessern,
  \item körperliche Fähigkeiten erweitern,
  \item besondere Bewegungsarten verleihen,
  \item außergewöhnliche Sinne ermöglichen,
  \item soziale Fähigkeiten stärken oder
  \item besonderes Wissen und handwerkliche Kenntnisse vermitteln.
\end{itemize}

Dadurch können sich zwei Helden mit derselben Rolle völlig unterschiedlich entwickeln.

\section{Trainierte und untrainierte Talente}

Ein Held kann im Laufe seiner Abenteuer viele Talente erlernen.

Gleichzeitig kann er jedoch nur einen Teil davon aktiv nutzen.

Ein Held kann gleichzeitig bis zu

\begin{itemize}
  \item vier Rollentalente und
  \item vier allgemeine Talente
\end{itemize}

trainiert haben.

Nur trainierte Talente können eingesetzt werden.

Auch nicht trainierte Talente gelten weiterhin als erlernt. Sie können deshalb Voraussetzungen für andere Talente erfüllen.

Wie ein Held neue Talente erlernt, wird in Kapitel 9 -- Wie entsteht ein Held? beschrieben.

\section{Talente neu trainieren}

Im Laufe seiner Abenteuer kann ein Held lernen, andere Talente einzusetzen.

Erhält ein Held die Gelegenheit, neu zu trainieren, darf er beliebig viele seiner trainierten Talente austauschen.

Dabei gelten folgende Regeln:

\begin{itemize}
  \item Es können höchstens vier Rollentalente trainiert sein.
  \item Es können höchstens vier allgemeine Talente trainiert sein.
  \item Rollentalente können nur gegen andere bereits erlernte Rollentalente ausgetauscht werden.
  \item Allgemeine Talente können nur gegen andere bereits erlernte allgemeine Talente ausgetauscht werden.
\end{itemize}

Nach dem Training können die neu gewählten Talente sofort eingesetzt werden.

\section{Wie werden Talente eingesetzt?}

Wie ein Talent eingesetzt wird, steht auf seiner Talentkarte.

Manche Talente wirken dauerhaft. Solange ein Talent trainiert ist, gilt der beschriebene Effekt jederzeit.

Andere Talente werden nur in bestimmten Situationen ausgelöst. Dafür kann zum Beispiel ein Angriff, eine Probe, ein Treffer oder eine andere Handlung notwendig sein.

Manche Talente verlangen außerdem, dass der Held:

\begin{itemize}
  \item einen Mondkristall ausgibt,
  \item die Talentkarte umdreht,
  \item eine bestimmte Handlung verwendet oder
  \item eine andere auf der Talentkarte genannte Bedingung erfüllt.
\end{itemize}

Wird eine Talentkarte umgedreht, ist sie erschöpft. Sie kann erst wieder verwendet werden, nachdem sie bereitgemacht wurde.

\subsection{Mehrere Talente gleichzeitig einsetzen}

Ein Held kann mehrere trainierte Talente besitzen, die in derselben Situation eingesetzt werden können.

Solange die Bedingungen aller Talente erfüllt sind, können ihre Wirkungen miteinander kombiniert werden.

Dabei wird jedes Talent einzeln ausgeführt. Kosten müssen für jedes Talent getrennt bezahlt werden.

Wenn sich zwei Talente widersprechen, gilt die speziellere Regel der Talentkarte.

\section{Der eigene Heldenpfad}

Talente bestimmen den Weg, den ein Held im Laufe seiner Abenteuer einschlägt.

Rollentalente legen fest, welche Aufgabe ein Held innerhalb seiner Gruppe übernimmt. Allgemeine Talente verleihen ihm darüber hinaus besondere Fähigkeiten und Eigenschaften.

Dadurch können zwei Helden mit ähnlichen Attributen, Fertigkeiten und derselben Ausrüstung dennoch vollkommen unterschiedlich sein.

Ein Held kann sich auf eine einzelne Rolle spezialisieren oder Rollentalente aus mehreren Rollen miteinander verbinden. Gleichzeitig kann er durch allgemeine Talente besondere Fähigkeiten entwickeln, die unabhängig von seiner Rolle sind.

Es gibt dabei keinen vorgeschriebenen richtigen Weg. Die Talente zeigen, welche Erfahrungen ein Held gesammelt hat und zu welcher Art von Held er im Laufe seiner Abenteuer geworden ist.